%% Preface — 그 위의 계층
\frontchapter{그 위의 계층}{서문}
\chaptersubtitle{운영체제 교수가 운영체제 위의 세계에 관한 책을 쓴 까닭}

나는 17년 동안 학부 운영체제를 가르쳤고, 매번 같은 고백으로 수업을 시작했다. \textbf{여러분 가운데 운영체제를 직접 만들 사람은 거의 없다.} 세상은 손에 꼽을 만큼 적은 커널 위에서 돌아가며, 그 대부분은 강의실의 학생들보다 나이가 많다. 새 커널의 시장 규모는 아무리 후하게 반올림해도 0에 가깝다. 그러면 해마다 용감한 학생 하나가 손을 든다. \textit{그런데 우리는 왜 이 과목을 배우나요?} 이 수업은 애초부터 운영체제라는 결과물 자체를 배우는 자리가 아니었기 때문이다. 그 결과물을 이루는 생각, 곧 스케줄링, 캐싱, 가상 메모리, 격리, 그리고 천 개의 프로세스에 동시에 약속하고도 그중 충분히 많은 약속을 지키는 기술을 배우는 자리였다. 또 학생이 컨텍스트 스위칭은 마법이 아니라 터무니없이 빠른 장부 정리일 뿐임을 깨닫는 바로 그 순간을 지켜보는 자리이기도 했다. 학생은 지적 불꽃을 얻고, 나는 그 불꽃이 튀는 모습을 보았다. 더없이 만족스러운 거래였다.

그런데 땅이 움직였다. 앞으로 학생들이 코드를 직접 쓰는 양이 갈수록 줄어들고, 편집기에 자료구조를 타이핑하는 일은 스스로 \textit{하는} 일보다 누군가에게 \textit{지시하는} 일이 되리라는 전망은 이제 대담한 주장도 아니다. 여전히 프로그래밍 솜씨의 가치가 예전과 다르지 않을 것이라고 말하는 동료도 있다. 나는 그만큼 낭만적이지 않다. 계산기가 나온 뒤에도 산술은 여전히 진짜 기술이고 때로는 결정적이지만, 더는 한 사람의 경력을 다른 사람과 가르는 기준이 아니다. 프로그래밍도 그렇게 될 것이다. 나는 이 말을 기쁜 마음으로 하는 것이 아니다. 나 역시 그 솜씨를 좋아했다. 하지만 머리말이라면 적어도 저자의 진짜 생각만큼은 솔직히 밝혀야 한다.

\subsection*{다시, 1968년}

땅이 움직일 때는 이전에 어디에서 움직였는지 돌아보면 마음을 다잡을 수 있다. 1968년 가을, 뜻밖에도 NATO가 바이에른 가르미슈에 약 50명을 모았다. 일부러 도발적으로 고른 기치가 \textit{소프트웨어 공학(software engineering)}이었다. 당시 많은 사람에게 두 단어의 조합 자체가 모순처럼 들렸다. 모임의 계기는 모두가 사석에서 \textbf{소프트웨어 위기}라고 부르던 상황이었다. 곳곳의 프로젝트가 일정을 넘기고, 예산을 초과하고, 틀린 결과물을 내놓았다. 소프트웨어가 독립된 학문 분야로 태어난 날을 하나 고르라면 바로 그 주만큼 어울리는 때도 없다.

아홉 달 뒤, 그 주장의 의미는 높은 하늘 위에서 입증되었다. 아폴로 11호가 달로 하강하던 중 잘못 설정된 레이더가 무의미한 작업을 쏟아 내자 유도 컴퓨터에 1202와 1201 경보가 번쩍이기 시작했다. 마거릿 해밀턴이 이끈 MIT 팀의 소프트웨어는 우선순위가 낮은 작업을 버리고 착륙 계산을 계속 수행했다. 언론의 제목을 장식한 것은 경보 코드였지만, 착륙을 구한 것은 \textbf{스케줄링 정책}이었다. 그보다 한 달 전 IBM은 소프트웨어를 기계와 분리해 판매하는 이른바 "언번들링"을 발표했고, 소프트웨어는 제품이자 산업이며 돈이 되었다. 보이지 않는 시스템 소프트웨어가 눈앞의 결과를 결정했다. 1960년대를 한 문장으로 요약하면 그렇다.

나는 지금 그 역사가 한 층 위에서 재판되고 있다고 믿는다. 오늘날의 AI 서비스는 운영체제 \textit{위}에 놓인, 그에 못지않게 보이지 않는 시스템 소프트웨어 계층에 서 있다. 하이퍼바이저는 실리콘을 조각내 빌려 주고, 컨테이너 런타임과 쿠버네티스는 수많은 머신을 한 대의 컴퓨터처럼 바꾸며, 서빙 엔진은 GPU 메모리에 천 개의 대화를 밀어 넣는다. 그 위로는 에이전트형 AI의 하니스와 파이프라인이 빠르게 늘어난다. 무대에서 이 계층 자체를 시연하는 사람은 없지만, 무대 위 모든 시연이 이 계층에 기대고 있다.

운영체제 교수에게 가장 즐거운 대목은 이 모든 것이 낯익다는 사실이다. 연구자들이 LLM \textbf{KV 캐시}의 메모리 낭비를 마침내 길들였을 때, 해법이 된 PagedAttention의 핵심은 페이지 테이블이었다. 1962년 Atlas 머신이 도입한 바로 그 기술이다. 제11장의 한 절 제목은 "운영체제가 안으로 들어오다"인데, 이는 비유가 아니다. 스케줄링은 연속 배칭으로, 격리는 네임스페이스와 cgroup으로 돌아왔다. 대기행렬 이론은 p99 지연 시간이라는 이름표를 달고 다시 나타났다. 내가 17년 동안 가르친 내용은 하나도 버려지지 않았다. 이름을 바꾸고 GPU 위에서 실행되는 모습으로 모두 돌아왔다.

\subsection*{두 번 가르치고 폐강한 과목}

이 책에는 실패한 선조가 있다. 독자에게도 그 집안 내력을 들려줄 필요가 있다. 챗봇 덕분에 "토큰"이라는 말이 일상어가 되기 전, 나는 \textit{클라우드와 AI}라는 대학원 과목을 두 차례 가르쳤다. 방향은 옳았지만 변화의 속도를 감당하지 못한 과목이었다. 쿠버네티스가 아직 전쟁에서 승리하기 전이어서, 진지한 사람들도 Mesos와 Swarm에 베팅했다. 9월에 가르친 AI 프레임워크가 12월이면 구식이 되었다. 첫 주에 쓴 강의계획서는 15주 차가 되기도 전에 고쳐야 했다. 내가 원한 교재, 곧 릴리스 노트보다 오래 살아남을 원리를 중심으로 구성한 책은 세상에 없었다. 결국 여기저기서 빌린 슬라이드를 희망으로 간신히 이어 붙여 강의했다.

학생들이 그 과목을 즐겼다고는 생각하지 않는다. 적어도 나는 분명 즐기지 못했다. 두 번째 강의를 마친 뒤 과목을 폐강했다. 더 나은 한 학기를 누릴 자격이 있는 주제를 제대로 가르치지 못했다는 교사 특유의 죄책감이 남았다. 이 책은 무엇보다도 그 빚을 갚는 일이다.

\subsection*{왜 지금인가}

달라진 점은 이제 뼈대가 굳었다는 것이다. 트랜스포머는 아홉 살이 되었고, 이를 송두리째 대체하려던 시도는 아직 왕좌를 빼앗지 못했다. 도전자들은 트랜스포머 곁에서 하이브리드 형태로 살아남았고, 가장 값비싼 부작용인 KV 캐시는 여전히 이 분야의 대표적인 미해결 과제로 남아 이제는 그 자체로 가르칠 만한 주제가 되었다. 서빙 스택도 연속 배칭, 대화를 위한 페이지드 메모리, 지연 시간을 재는 두 시계라는 알아볼 수 있는 형태로 수렴했다. 쿠버네티스는 전쟁에서 완승한 나머지 업계 언론이 이제 \textit{지루하다}고 부른다. 인프라 분야에서 받을 수 있는 최고의 찬사다. 클라우드 네이티브 아키텍처는 조직이 시스템을 만드는 보편적인 방식이 되었다. 앞으로도 하니스, 런타임, 파이프라인 같은 새 모듈은 이 뼈대 위에 계속 등장하겠지만, 모두 이 뼈대 \textit{안으로} 들어올 것이다. 세부 사항이 매주 바뀔 때는 블로그 글을 쓴다. 원리가 10년은 버틸 듯할 때 비로소 조심스럽게 교재를 쓸 수 있다.

\subsection*{이 책을 만든 방법: 공개해야 할 이야기}

교재가 필요하다는 사실을 아는 것과 실제로 쓰는 것은 전혀 다른 종목이다. 도널드 커누스는 1962년 컴파일러에 관한 소박한 책 한 권을 쓰기로 계약했다. 첫 권은 6년 뒤에 나왔고, 60여 년이 지난 지금도 그 장대한 프로젝트는 미완성이다. 프레드 브룩스는 프로젝트가 어떻게 1년씩 늦어지는지 묻고 스스로 답했다. \textit{하루씩 늦어진 결과다.} 교재도 똑같다. 그렇다고 대학원생 한 명에게 한 장씩 맡긴 뒤 교수 이름으로 묶어 내던 학계의 전통적인 해결책을 꺼낼 수도 없었다. 그런 시대는 끝났고, 마땅히 그래야 한다. 대학원생이 세상에 내놓아야 할 것은 연구이지 내 책의 원고가 아니다.

이 책을 구한 것은 다름 아닌 이 책의 주제였다. 나는 AI 에이전트를 공동 저자로 삼는 일에 회의적인 태도로 시작했다. 결과를 있는 그대로 말하자면, 이 책을 손으로 썼다면 지금보다 더 잘 쓰지 못했을 것이며 이번 10년 안에 끝내지도 못했을 것이다. 내가 보탠 것은 뼈대였다. 강의, 강의계획서, 무엇이 근본 원리이고 무엇이 유행인지 가리는 판단이었다. 에이전트는 초안과 그림, 계산, 그리고 지칠 줄 모르는 인내를 보탰다. 나는 동그라미 하나와 선 다섯 개로 사람을 그린 연필 스케치를 건넸고, 영화 같은 결과물을 돌려받았다. 농기계가 한 쌍의 손으로 경작할 수 있는 면적을 몇 배로 늘렸듯, 이제 지적 노동에도 같은 배율의 증폭이 찾아왔다. (전통적으로 머리말은 경건한 문구로 끝난다. 도움을 준 이들에게 감사를 돌리고, 남은 오류는 오롯이 저자의 몫이라고 말한다. 이 시대에 맞게 고쳐 보겠다. 이 책에서 좋은 부분은 모두 내 덕이고, 남은 오류는 전적으로 AI 탓이다. 왜냐하면 사람은 여전히 중요하기 때문이다.)

공개할 이야기 안에 한 가지가 더 있다. 초안의 모든 문장은 빌린 GPU에서 실행되는 모델이 토큰으로 만들어 내게 보낸 것이다. 그 아래에서는 요청이 연속으로 배칭되고, 대화가 KV 캐시를 통해 페이지 단위로 오갔으며, 제9장부터 제12장까지 다루는 장치 전체가 쉼 없이 돌아갔다. 이 책은 자신이 가르치는 바로 그 시스템으로 만들어졌다. 이보다 강의계획서를 더 강하게 뒷받침하는 증거는 떠올리기 어렵다.

\subsection*{이 책의 내용과 읽는 방법}

이 책은 15주 강의를 따라 세 부로 구성된다. \textbf{1부: 인프라와 원리(제1장\textasciitilde{}제6장)}에서는 의도적으로 스택을 내려간다. 클라우드 인스턴스와 컨테이너의 실체(네임스페이스, cgroup, 계층형 파일 시스템이라는 세 커널 기능이 트렌치코트를 걸친 모습), 쿠버네티스가 선언된 의도를 실제로 실행되는 머신 집합으로 바꾸는 방법, 가중치와 데이터가 머무는 곳을 살펴본다. \textbf{2부: AI 서빙 시스템(제7장, 제9장\textasciitilde{}제12장)}에서는 꼬리 지연 시간과 SLO, 서빙 프레임워크, 양자화와 컴파일, LLM 추론만의 물리 법칙, 비용표를 붙인 스케일링을 통해 서빙이라는 분야의 기초를 세운다. \textbf{3부: 운영(제13장\textasciitilde{}제14장)}에서는 레지스트리, 파이프라인, 카나리를 이용해 소란 없이 시스템을 바꾸는 법과 모니터링, 알림, 드리프트로 이루어진 서비스의 신경계를 배운다. 제8장은 없다. 중간고사에서 독자가 직접 쓰는 장이기 때문이다. 마지막에는 맺음말이 이 모든 이야기를 문밖까지 배웅한다.

전체 내용은 하나의 연속된 이야기로 이어진다. 학생 팀의 데모는 21시 04분, 정확히 100명의 사용자가 몰리자 무너진다. 이를 "고치기" 위해 빌린 GPU는 대부분 놀았지만 월말 청구서는 779.66달러다. 15주 동안 엔지니어링을 거친 뒤 같은 서비스는 약속한 범위 안에서 월 900달러에 운영된다. \textbf{정확히} 900달러다. 그때는 비용을 당하는 것이 아니라 스스로 선택하기 때문이다. 각 장은 새벽 2시에도 결정을 내릴 수 있을 만큼 선명한 한 줄짜리 \textit{기념품}을 남긴다. \code{exit 137}은 커널이 남기는 작별 인사이고, \textit{평균은 아무도 설명하지 못한다.} 풀이 예제는 독자 앞에서 계산 과정을 숨김없이 보여 준다. 좋은 세미나라면 질문을 던질 바로 그 지점에서 \textit{생각해 보기} 상자가 흐름을 멈춘다. 모든 연습문제에는 정답과 해설이 딸려 있다. 이 책은 \textbf{혼자} 공부하는 사람도 읽을 수 있어야 하기 때문이다.

이는 의도한 구성이다. 지금은 AI 시대다. 설명은 값싸졌다. 새벽 3시에도 AI 튜터가 무엇이든 기꺼이 설명해 준다. 하지만 \textbf{큐레이션은 여전히 값싸지 않다.} 무엇을 어떤 순서로 배우고, 어떤 원리로 연결할 것인가. 그것이 강의이며, 이 책은 종이 위의 강의가 되려 한다. 학생에게 뼈대와 질문을 건네면 나머지는 학생과 AI가 해낼 것이다. (동료 교수들에게: 사직서는 아니다. 여전히 누군가는 뼈대를 골라야 한다.) 형식적인 선수 요건은 다음과 같다. Python은 필수다. 운영체제 과목은 강력히 권한다. 필요한 내용은 책에서 다시 세우므로 반드시 이수해야 하는 것은 아니다. 기계학습 입문 과목도 도움이 되지만 전제로 삼지는 않는다.

\vspace{1.0pt}

\begin{calloutbox}{tbbamber}{tbbamberfill}
{\normalsize\bfseries\color{tbbamberdark}이 책을 활용하는 방법}\par\vspace{3pt}
\begin{tbbitemize}
  \item \textbf{정규 강의}: 각 장은 15주 학기의 한 주와 일대일로 대응한다. 8주 차에는 중간고사, 15주 차에는 데모 데이가 있으며, 별도 강의계획서에 프로젝트 평가 기준과 실습 일정을 담았다.
  \item \textbf{독학}: 한 장을 읽고 실습한 뒤 AI 튜터와 연습문제의 답을 놓고 논쟁해 보자. 모든 문제에 정답과 해설이 있어 논쟁의 승자를 스스로 판정할 수 있다.
  \item \textbf{하드웨어}: 모든 실습은 소형 GPU 한 개면 충분하다. 무료 Colab T4로도 가능하며 CPU만 사용하는 대안도 제공한다. 보유하지 않은 클러스터가 필요한 실습은 없다.
\end{tbbitemize}
\end{calloutbox}

\vspace{1.5pt}

\subsection*{"시스템SW가 돈이다"}

얼마 전 한 대형 기술 기업의 최고경영자가 이 모든 내용을 네 단어로 요약했다. 이니셜은 S다. 한국 독자라면 누구인지 맞히는 재미가 있을 것이다. \textbf{"시스템SW가 돈이다."} 1969년 IBM도 같은 결론에 도달했으니 훌륭한 선배가 있는 셈이다. 이 글을 쓰는 현재 세계에서 가장 가치 있는 기업은 명함상 반도체 회사다. 하지만 그 플랫폼을 떠나 보려 한 사람에게 물으면 진짜 해자가 CUDA임을 알게 된다. 시스템 소프트웨어는 언제나처럼 보이지 않는 곳에서 조용히 대금을 거둔다.

독자에게 바라는 것은 두 가지다. 첫째, 이 책을 다 읽은 뒤에는 보이지 않는 계층을 \textit{볼 수} 있기를 바란다. 챗봇이 0.5초 만에 답하는 모습을 보면서 그 뒤의 하이퍼바이저, cgroup, 스케줄러, 페이지드 캐시, 청구서를 함께 떠올릴 수 있기를 바란다. 둘째, 여러분 세대의 21시 04분이 찾아왔을 때, 그리고 그 순간은 반드시 찾아온다, 그 방 안에서 어느 구간이 병목인지 이미 아는 사람이 바로 여러분이기를 바란다. 1960년대에는 그런 시야가 수많은 경력을 만들어 냈다. 2020년대에는 한 층 위에서 같은 일이 다시 벌어지고 있다.

아직 답이 없던 시절 이 질문들을 베타 테스트해 준 폐강 과목의 두 기수 학생들에게 감사한다. 내가 졸라맨을 그리는 동안 쉬지 않고 타이핑한 에이전트들에게도 감사를 전한다.

\textcolor{tbbink}{\textbf{여러분의 서비스가 지루할 만큼 평온하기를.}}

{\raggedleft \textcolor{tbbink}{\textbf{의성}}\textcolor{tbbgray}{\textit{  ·  한국 AI 혁명의 심장 판교에서, 2026년 7월}}\par}
